\documentclass[10pt,letterpaper]{article}
\usepackage[T1]{fontenc}
\usepackage{amssymb}
\usepackage{amsmath}
\usepackage[margin=0.5in]{geometry}
\usepackage{siunitx}
\usepackage{xcolor}
\title{Introduction to Quantum Electronic Devices Post Lecture 16 Exercise Problem 1}
\author{Jeremy Chen}
\date{23/03/2026}
\begin{document}
	\maketitle

\noindent In metals, $E_{F}$ is the energy of the highest-energy electron in the metal, therefore $E$ for all electrons satisfies the following constraint: $E \leq E_{F}$.

\noindent What is the occupation probability of a level whose energy is 1 eV below $E_{F}$ (i.e. $E - E_{F} = -1$ eV)?

\color{violet}
Given the Fermi-Dirac distribution, the occupation probability is
\begin{gather*}
f(E) = \frac{1}{1 + \exp{\left( \frac{E - E_{F}}{K_{B}T} \right )}} = \frac{1}{1 + \exp{\left( \frac{-1}{300 * \num{8.617e-5}} \right)}} = 1
\end{gather*}
\color{black}

\end{document}